\appendix
\section{Additional Experimental Results}
\subsection{Logit Lens KL Divergence}
\label{appA-subsection:kl-divergence}

Another visualization of the logit lens represents the KL divergence between logits produced
at every layer and the final layer. In Figure~\ref{fig:logit-lens-KL} darker colors suggest a larger KL divergence
between the distributions, meaning that the logits distributed very different compared the the outer layer one.

\begin{figure}[H]
    \centering
    \includegraphics[width=\textwidth]{hidden_vs_outpout_KL.png}
    \caption{logit lens: KL divergence between each hidden layer output and the final output}
    \label{fig:logit-lens-KL}
\end{figure}

\noindent An interesting observation in Figure~\ref{fig:logit-lens-KL} concerns the first token. Its KL divergence remains 
consistently high throughout most of the network and drops sharply only in the final two layers. 
This suggests that the representation of the first token goes through a more aggressive refinement near the end of the 
computation, with the final layers contributing considerably more to the output distribution.


\subsection{Memory State Norm tracking}
\label{appA:memory-norm}

We extend the input and forget gate analysis by examining how much information is stored in the memory state.
To quantify this, we compute the Frobenius norm of the cell state$C_t \in \mathbb{R}^{d \times d}$.
As in Subsection~\ref{subsub: gate-dynamics}, we track the contribution of each token to the evolution of the state norm,
where color indicates the norm after processing each token. 
A time series visualization is shown in Figure~\ref{fig:frobenius time series}.

\begin{figure}[H]
    \centering
    \includegraphics[width=\textwidth]{frobenius_c_t_colored.png}
    \caption{Individual influence of tokens to the norm of the memory state}
    \label{fig:frobenius_colored}
\end{figure}
\begin{figure}[H]
    \centering
    \includegraphics[width=\textwidth]{c_t_norm_time_series.png}
    \caption{Time series visualization of Figure~\ref{fig:frobenius_colored}}
    \label{fig:frobenius time series}
\end{figure}

\noindent The visualization shows that the memory state gradually builds up information over time, as indicated by the 
steady increase in the Frobenius norm along the sequence. We also observe sharp spikes in magnitude when the 
model processes information-rich proper nouns, suggesting these tokens require stronger updates to encode 
their meaning in the state, while also noticing a drop in the norm value when processing punctuation tokens. 

\section{SQuAD dataset sampling and evaluation details}
\label{app-sec:squad-details}
\subsection{Prompt}
\begin{tcolorbox}[
    colback=gray!5!white,
    colframe=blue!70!black,
    title=\textbf{SQuAD Evaluation Prompt Template},
    fonttitle=\sffamily\bfseries,
    boxrule=1mm,
    arc=3mm
]
\noindent\textbf{Context:}\\[1em]
\textit{\{context\_str\}}\\[1em]
\textbf{Question:} \textit{\{question\}}\\
\textbf{Answer:} 
\end{tcolorbox}

\subsection{SQuAD Model Evaluation Examples}
\label{app:squad-examples}

% Note: To compile this code, ensure you have \usepackage{tcolorbox} in your document preamble.

The following examples present randomly selected questions from our curated SQuAD dataset alongside the answers generated by xLSTM (Cache), Mistral, and xLSTM (No Cache). Context has been truncated. 

% ---------------------------------------------------------
% EXAMPLE 1
% ---------------------------------------------------------
\subsubsection{Example 1: Philosophy \& Physics}

\begin{tcolorbox}[colback=blue!5!white, colframe=blue!75!black, title=Prompt]
\textbf{Context:} Philosophers in antiquity used the concept of force in the study of stationary and moving objects and simple machines, 
but thinkers such as Aristotle and Archimedes retained fundamental errors in... (text continues)\\
\textbf{Question:} Who had mathmatical insite?
\end{tcolorbox}

\begin{tcolorbox}[colback=gray!5!white, colframe=gray!75!black, title=\textbf{xLSTM (Cache)}, ]
\begin{verbatim}
Philosophers in antiquity used the concept of force in the study 
of stationary and moving objects and simple machines, but
\end{verbatim}
\end{tcolorbox}

\begin{tcolorbox}[colback=gray!5!white, colframe=gray!75!black, title=\textbf{Mistral}]
\begin{verbatim}
Sir Isaac Newton

Question: What was the theory of relativity?
Answer: A theory of relativity that
\end{verbatim}
\end{tcolorbox}

\begin{tcolorbox}[colback=gray!5!white, colframe=gray!75!black, title=\textbf{xLSTM (No Cache)}]
\begin{verbatim}
Philosophers in antiquity used the concept of force in the study 
of stationary and moving objects and simple machines, but
\end{verbatim}
\end{tcolorbox}

\vspace{1em}

% ---------------------------------------------------------
% EXAMPLE 2
% ---------------------------------------------------------
\subsubsection{Example 2: History of Civil Disobedience}

\begin{tcolorbox}[colback=blue!5!white, colframe=blue!75!black, title=Prompt]
\textbf{Context:} One of its earliest massive implementations was brought about by Egyptians against 
the British occupation in the 1919 Revolution. Civil disobedience is one of the many ways people have rebelled... (text continues)\\
\textbf{Question:} Where was the location of the 2003 Rose revolution?
\end{tcolorbox}

\begin{tcolorbox}[colback=gray!5!white, colframe=gray!75!black, title=\textbf{xLSTM (Cache)}]
\begin{verbatim}
Answer:

One of its earliest implementations was brought about by the 
British in the fight against the British Civil War.
\end{verbatim}
\end{tcolorbox}

\begin{tcolorbox}[colback=gray!5!white, colframe=gray!75!black, title=\textbf{Mistral}]
\begin{verbatim}
Georgia
\end{verbatim}
\end{tcolorbox}

\begin{tcolorbox}[colback=gray!5!white, colframe=gray!75!black, title=\textbf{xLSTM (No Cache)}]
\begin{verbatim}
Answer:

One of its earliest implementations was brought about by the 
British in the fight against the British Civil War.
\end{verbatim}
\end{tcolorbox}

\vspace{1em}

% ---------------------------------------------------------
% EXAMPLE 3
% ---------------------------------------------------------
\subsubsection{Example 3: European Geography (City)}

\begin{tcolorbox}[colback=blue!5!white, colframe=blue!75!black, title=Prompt]
\textbf{Context:} The Rhine (Romansh: Rein, German: Rhein, French: le Rhin, Dutch: Rijn) is a European 
river that begins in the Swiss canton of Graubünden in the southeastern Swiss Alps, forms part... (text continues)\\
\textbf{Question:} What is the largest city the Rhine runs through? 
\end{tcolorbox}

\begin{tcolorbox}[colback=gray!5!white, colframe=gray!75!black, title=\textbf{xLSTM (Cache)}]
\begin{verbatim}
The Rhine runs through the largest city the Rhine runs through, 
the Rhine, the Rhine, the
\end{verbatim}
\end{tcolorbox}

\begin{tcolorbox}[colback=gray!5!white, colframe=gray!75!black, title=\textbf{Mistral}]
\begin{verbatim}
Cologne, Germany

Question: What is the second largest city the Rhine runs through?
Answer:
\end{verbatim}
\end{tcolorbox}

\begin{tcolorbox}[colback=gray!5!white, colframe=gray!75!black, title=\textbf{xLSTM (No Cache)}]
\begin{verbatim}
The Rhine runs through the largest city the Rhine runs through, 
the Rhine, the Rhine, the
\end{verbatim}
\end{tcolorbox}

\vspace{1em}

% ---------------------------------------------------------
% EXAMPLE 4
% ---------------------------------------------------------
\subsubsection{Example 4: European Geography (River)}

\begin{tcolorbox}[colback=blue!5!white, colframe=blue!75!black, title=Prompt]
\textbf{Context:} The Rhine (Romansh: Rein, German: Rhein, French: le Rhin, Dutch: Rijn) is a European river 
that begins in the Swiss canton of Graubünden in the southeastern Swiss Alps, forms part... (text continues)\\
\textbf{Question:} What river is larger than the Rhine?
\end{tcolorbox}

\begin{tcolorbox}[colback=gray!5!white, colframe=gray!75!black, title=\textbf{xLSTM (Cache)}]
\begin{verbatim}
The Danube (Romansh: Danub, German: Donau, French: le Danube, Dutch
\end{verbatim}
\end{tcolorbox}

\begin{tcolorbox}[colback=gray!5!white, colframe=gray!75!black, title=\textbf{Mistral}]
\begin{verbatim}
The Danube is larger than the Rhine.
\end{verbatim}
\end{tcolorbox}

\begin{tcolorbox}[colback=gray!5!white, colframe=gray!75!black, title=\textbf{xLSTM (No Cache)}]
\begin{verbatim}
The Danube (Romansh: Danub, German: Donau, French: le Danube, Dutch
\end{verbatim}
\end{tcolorbox}

\vspace{1em}

% ---------------------------------------------------------
% EXAMPLE 5
% ---------------------------------------------------------
\subsubsection{Example 5: Sociology \& Law}

\begin{tcolorbox}[colback=blue!5!white, colframe=blue!75!black, title=Prompt]
\textbf{Context:} One of its earliest massive implementations was brought about by Egyptians against the British 
occupation in the 1919 Revolution. Civil disobedience is one of the many ways people have rebelled... (text continues)\\
\textbf{Question:} What is it called when people in society rebel against laws they think are unfair?
\end{tcolorbox}

\begin{tcolorbox}[colback=gray!5!white, colframe=gray!75!black, title=\textbf{xLSTM (Cache)}]
\begin{verbatim}
Answer:

One of its earliest implementations was brought about by people 
in society rebelling against the laws they deemed unfair
\end{verbatim}
\end{tcolorbox}

\begin{tcolorbox}[colback=gray!5!white, colframe=gray!75!black, title=\textbf{Mistral}]
\begin{verbatim}
Civil disobedience.

Question: What is the difference between civil disobedience and 
civil disobed
\end{verbatim}
\end{tcolorbox}

\begin{tcolorbox}[colback=gray!5!white, colframe=gray!75!black, title=\textbf{xLSTM (No Cache)}]
\begin{verbatim}
Answer:

One of its earliest implementations was brought about by people 
in society rebelling against the laws they deemed unfair
\end{verbatim}
\end{tcolorbox}